\subsection{Logiciels utilisés}

\noindent Les principaux logiciels utilisés sont :
\begin{itemize}
    \item \textbf{\gls{vscode}}, qui est un \gls{ide} (environnement de développement intégré) permettant de développer en Python, C++ et d'autres langages. Il est utilisé pour écrire le code des algorithmes de contrôle en Python. \gls{vscode} dispose de nombreuses extensions, dont une extension \textit{ssh} permettant de se connecter à distance à la Raspberry Pi des robots. Cela permet d'avoir accès aux fichiers de cette dernière via l'explorateur de fichiers de l'\gls{ide}, et de modifier les fichiers directement depuis l'\gls{ide}, facilitant grandement le développement.
    \item \textbf{Matlab/Simulink}, qui est un environnement de développement pour la simulation et la modélisation de systèmes dynamiques. Il a surtout été utilisé pour visualiser les données en temps réel, et tracer les différents graphiques donnés en partie \ref{chap:resultats}.
    \item \textbf{\gls{ide} Arduino}, afin de programmer la carte OpenCR des robots.
    \item \textbf{GitHub}, qui est une plateforme en ligne permettant d’héberger du code, de le versionner et de collaborer à plusieurs sur un même projet. GitHub est particulièrement utile pour les systèmes distribués, car il suffit d'utiliser la commande \textit{git clone} sur chaque robot pour récupérer la dernière version du code, ce qui simplifie grandement la synchronisation du projet sur plusieurs machines.
\end{itemize}

L'ensemble des programmes réalisés durant le stage est disponible sur le Github suivant : 
\newline \url{https://github.com/watson67/mecanum}

\noindent Différents fichiers \textit{README} s'y trouvent pour des explications plus détaillées. 
